\paragraph{MIP pool expansion}
\label{sec:mip_pool_expansion}
The constructed pool is feasible and cell-complete, but completeness is a
property of the pool, not of a zone: in a given zone a control-table cell
$(k,m)$ can require more records than that zone's slice of the pool can supply,
and in the limiting case no household in the slice carries the cell at all.
This is the zero-cell problem, and no weight vector or selection can repair it,
because the records needed do not exist. Expansion therefore treats the
zero-capacity case as the constraint it is, once the audit has established which
cells are in it.

\emph{What the audit measures.} For zone $z$ and control cell $(k,m)$, let
$E_{z,k,m}$ be the census target and $s_{z,k,m}$ the number of candidate records
in the zone's pool slice that carry the cell. The audit compares the target with
the \emph{attainable capacity} $r_{\max}s_{z,k,m}$, where $r_{\max}$ is the
largest number of times the zone fitter may replicate a single candidate, and
places every cell with $E_{z,k,m}>0$ in one of three states: \emph{sufficient}
if $r_{\max}s_{z,k,m}\ge E_{z,k,m}$; a \emph{partial shortfall} if
$0<r_{\max}s_{z,k,m}<E_{z,k,m}$, which replication can reach and which is
therefore left to the fitter; and \emph{zero capacity} if $s_{z,k,m}=0$, which
no weight, selection or replication can reach because the candidates are absent.
Only the third state is a pool deficiency, and only the third state is what
expansion addresses. The distinction matters for how a poor zone fit is read: a
cell in the second state is a fitting outcome, a cell in the third is a support
defect that no fitter could have avoided, and the two are reported separately
rather than pooled into one error term. \Cref{fig:mip_expansion} follows one zone
through the step, and Table~\ref{tab:cp_vs_mip} states how construction and
expansion differ.

\begin{figure}[t]
  \centering
  % =============================================================================
% fig:mip_expansion -- filling a zone's zero-capacity cells.
%
% Same board language as fig:cp_pool, one zone wide: a zone's control board has
% a positive target in a square its pool slice cannot supply at all, and no
% weight vector, selection or replication can reach it. Expansion builds
% donor x shell candidates for exactly those squares, one program per zone
% selects a handful, and the board is complete.
%
% Square states: blue = the pool slice covers the target; amber = covered but
% short; red = a positive target with no supply.
%
% Editing rules: no \\ inside a braced group; keep card text short, because the
% board is only ~2 cm tall and the cards sit beside it.
% =============================================================================
\providecolor{ForestGreen}{rgb}{0.0,0.258,0.145}
\begin{tikzpicture}[
  font=\footnotesize,
  >=latex,
  sq/.style={draw=black!50, minimum width=0.44cm, minimum height=0.42cm,
             inner sep=0pt, line width=0.4pt},
  have/.style={sq, fill=blue!16, draw=blue!45!black},
  short/.style={sq, fill=orange!30, draw=orange!70!black},
  zero/.style={sq, fill=red!25, draw=red!70!black, line width=0.8pt},
  filled/.style={sq, fill=ForestGreen!20, draw=ForestGreen!70!black, line width=0.8pt},
  hdr/.style={font=\scriptsize, text=black!60, inner sep=1pt},
  rowlab/.style={font=\scriptsize, anchor=east, text=black!85, inner sep=2pt},
  tag/.style={font=\scriptsize, text=black!75, align=center, inner sep=2pt},
  flow/.style={->, thick, draw=black!60},
  card/.style={draw=black!45, rounded corners=2pt, align=left, inner sep=3pt,
               fill=white},
  val/.style={font=\tiny, text=black!85, inner sep=0pt},
]

% ============================================================ board, zone z
\node[font=\scriptsize\bfseries, anchor=south] at (2.27,0.62)
  {zone $z$: \textnormal{target / in pool}};

\foreach \c [count=\j] in {1,...,6} {
  \node[hdr] at ({0.95+(\j-0.5)*0.44},0.34) {\c};
}

% 1 = covered, 0 = covered but short, 2 = no supply at all.
\foreach \r/\row/\name in {%
  0/{1,1,1,1,1,1}/SEXP,%
  1/{1,1,1,1,1,1}/HEAP,%
  2/{1,1,1,2,1,1}/MSTP,%
  3/{1,1,1,1,1,1}/OCCP,%
  4/{1,0,1,1,1,1}/VEHRD}%
{
  \node[rowlab] at (0.82,{-\r*0.44}) {\name};
  \foreach \d [count=\j] in \row {
    \pgfmathsetmacro{\x}{0.95+(\j-0.5)*0.44}
    \pgfmathsetmacro{\y}{-\r*0.44}
    \ifnum\d=2
      \node[zero, label={[val]center:8/0}] at (\x,\y) {};
    \else\ifnum\d=0
      \node[short, label={[val]center:9/3}] at (\x,\y) {};
    \else
      \node[have] at (\x,\y) {};
    \fi\fi
  }
}

% ======================================================== candidate for a cell
\draw[flow] (3.62,-1.00) -- (4.05,-1.00);

\node[card, anchor=north west] (cand) at (4.10,0.36) {\parbox{2.75cm}{%
  \scriptsize\raggedright
  \textbf{One candidate}\\[2pt]
  \textbf{donor}: a person with the square\\[2pt]
  \textbf{shell}: a real dwelling, a member in that role\\[2pt]
  \textbf{candidate}: donor in shell $=$ one household}};

% ================================================================ one program
\draw[flow] (7.02,-1.00) -- (7.45,-1.00);

\node[card, draw=blue!55!black, fill=blue!6, anchor=north west, name=prog]
  at (7.55,0.36)
  {\parbox{2.45cm}{\scriptsize\raggedright
  \textbf{One program}\\[2pt]
  covers every zero square at once\\[2pt]
  budget 64 households, 256 persons\\[2pt]
  no donor $\Rightarrow$ reported, not invented}};

% ==================================================================== settled
\draw[flow] (10.15,-1.00) -- (10.55,-1.00);

\node[font=\scriptsize\bfseries, anchor=south] at (11.92,0.62)
  {after expansion: \textnormal{2 households appended}};

\foreach \c [count=\j] in {1,...,6} {
  \node[hdr] at ({10.62+(\j-0.5)*0.44},0.34) {\c};
}

\foreach \r/\row/\name in {%
  0/{1,1,1,1,1,1}/SEXP,%
  1/{1,1,1,1,1,1}/HEAP,%
  2/{1,1,1,3,1,1}/MSTP,%
  3/{1,1,1,1,1,1}/OCCP,%
  4/{1,3,1,1,1,1}/VEHRD}%
{
  \foreach \d [count=\j] in \row {
    \pgfmathsetmacro{\x}{10.62+(\j-0.5)*0.44}
    \pgfmathsetmacro{\y}{-\r*0.44}
    \ifnum\d=3
      \node[filled, label={[val]center:\textbf{2}}] at (\x,\y) {};
    \else
      \node[have] at (\x,\y) {};
    \fi
  }
}

% ================================================================ audit line
% Anchored below the tallest card rather than at a fixed y, so the note cannot
% collide with the card text when the wording changes.
\node[anchor=north west, align=left, text width=13.2cm, font=\scriptsize,
      inner sep=1pt] at ($(prog.south west)+(-7.50,-0.32)$)
  {\textbf{What the audit found across ten zones.} 123 zero-capacity cells, 73 of
   them sourceable from donors, 70 covered with 14 households (18 persons) --- one
   or two households and at most three persons per zone. The 53 left uncovered are
   reported rather than invented: 40 sit in codes absent from the extract, 10 in a
   padding code, and 3 had no candidate that survived the checks. The zone fitter,
   not this step, decides how often an appended household is replicated.};

\end{tikzpicture}

  \caption{How a zone's zero-capacity cells are filled. For each zone the audit compares
  every positive-target cell with the supply in that zone's pool slice; on the board, a
  filled square is a cell the slice covers, an amber square one it covers only partly, and
  a red square one it cannot supply at all. A square with no supply cannot be reached by
  any weight vector, selection or replication. Each such cell
  receives a candidate: a real donor who carries the cell, written into a real shell
  dwelling from the same region that already has a member in the donor's role, so the
  dwelling and family terms stay observed rather than invented. One program per zone then
  selects a small set of candidates covering all zero cells at once, and the selected
  households are appended to the pool with fresh identifiers. Coverage is soft: a cell with
  no donor anywhere is reported by the audit rather than synthesised, and the zone fitter
  (IPU, SA or MIP), not this step, decides how often an appended household is replicated.}
  \label{fig:mip_expansion}
\end{figure}


\begin{table}[h!]
\centering
\caption{The two pool-side programs answer different questions and are not
substitutes. Construction makes the pool complete over the modelled attributes for
every zone at once; expansion repairs the cells a single zone's slice cannot supply.
Neither stage allocates records to zones, which is why both are followed by a fitter.}
\label{tab:cp_vs_mip}
\footnotesize
\setlength{\tabcolsep}{5pt}
\begin{tabular}{@{}>{\raggedright\arraybackslash}p{1.7cm}
                 >{\raggedright\arraybackslash}p{5.6cm}
                 >{\raggedright\arraybackslash}p{6.3cm}@{}}
\toprule
 & \textbf{CP pool construction} & \textbf{MIP pool expansion} \\
\midrule
Question & Can every rule-feasible cell be represented in the pool at all? &
Which cells can this zone's slice not supply, and which households can fill them? \\
Scope & The whole census, all zones at once, all 19 control tables &
One zone, and only the cells the audit places in the zero-capacity state \\
Decision & Which category each modelled attribute takes, given a household shell and the
observed conditional distributions & Which candidates enter a small added set, one binary
per candidate, against all marginal tables at once \\
Guarantee & Every emitted record satisfies $\mathcal{R}$, and the pool is cell-complete
over the modelled attributes & Every added household satisfies $\mathcal{R}$, and every
sourceable zero cell receives at least one candidate \\
Does not & Use zone targets, match any zone total, or decide replication & Close the
zone's full deficit, fit its totals, or decide how often a candidate is used \\
Cost & 27.1 minutes and 5.0\,GB for the whole pool, once & About 100 seconds per zone,
and only where zero cells exist \\
Result here & 1{,}404{,}801 persons, no rule violations & 14 households (18 persons)
across the ten audit zones \\
\bottomrule
\end{tabular}
\end{table}

\emph{Candidate construction.} For a zone $z$ let $\mathcal{Z}_z^0$ be the set of
control cells with a positive census target and zero attainable capacity in the
zone's candidate pool. For each such cell the module draws donor records from the
census microdata that carry the cell, pairs each donor member profile with a real
shell household from the same \texttt{REGUCP} that has a member in the donor's
\texttt{RLHP} role, and writes the donor profile into that slot. If the donor
profile conflicts with the shell's dwelling terms or with the rule set
$\mathcal{R}$, a small completion program finds the closest feasible profile,
\begin{equation}
\min_{y}\;\sum_{a}\bigl(1-y_{a,\tau(a)}\bigr)
\quad\text{s.t.}\quad
\sum_{c} y_{a,c}=1,\;\;
y_{a_1,c_1}+y_{a_2,c_2}\le 1 \;\;\forall\,(a_1,c_1,a_2,c_2)\in\mathcal{C},
\label{eq:mip_repair}
\end{equation}
where $\tau(a)$ is the donor's code for attribute $a$ and $\mathcal{C}$ is the
set of forbidden category pairs obtained by evaluating the shared rule set over
every attribute domain mesh (2{,}027 cuts in the ten-zone development audit).
Equation~\eqref{eq:mip_repair} is a \emph{candidate repair}, not the pool
objective: it preserves a real donor profile wherever the hard rules allow, and
it is solved separately for each candidate before any selection. The
household-level fields are fixed to the shell, so every candidate remains atomic
and family-coherent; candidates that fail the shared
rule checks are discarded, never patched. For attributes outside
the person-level model (e.g.\ \texttt{tend}, \texttt{typp}, \texttt{indp},
\texttt{stup}, \texttt{assnp}), a person-level repair is not available, so the
module falls back to a whole real donor household copied verbatim; the observed
household is feasible by construction.

\emph{Joint-marginal selection.} The candidates are then selected together, not
cell by cell. Let $\mathcal{J}_z$ index the candidate households,
$x_j\in\{0,1\}$ indicate selection, and $s_j$ be the number of persons in
candidate $j$. For a control cell $(k,m)$, let $a_{j,k,m}$ be the candidate's
contribution---the number of its members in the cell for a person-level table, or
one if the household sits in the cell for a household-level table---and let
$u_{j,k}=s_j$ for a person-level table and $u_{j,k}=1$ for a household-level
table. With the zone target $E_{z,k,m}$ and target share
\begin{equation}
q_{k,m}=\frac{E_{z,k,m}}{\sum_{m'\in\mathcal{M}_k}E_{z,k,m'}},
\label{eq:mip_share}
\end{equation}
the selection program is
\begin{equation}
\begin{aligned}
\min_{x,e,d}\quad &
\lambda_0\sum_{c\in\mathcal{Z}_z^0} e_c
\;+\;\sum_{k\in\mathcal{K}}\sum_{m\in\mathcal{M}_k} w_{k,m}\,d_{k,m}
\;+\;\lambda_1\sum_{j\in\mathcal{J}_z} x_j \\
\text{s.t.}\quad &
\sum_{j\in\mathcal{J}_z} a_{j,c}\,x_j+e_c\ge 1,
&& c\in\mathcal{Z}_z^0,\\
& d_{k,m}\ge \sum_{j\in\mathcal{J}_z}
  \bigl(a_{j,k,m}-q_{k,m}u_{j,k}\bigr)x_j,
&& k\in\mathcal{K},\,m\in\mathcal{M}_k,\\
& d_{k,m}\ge -\sum_{j\in\mathcal{J}_z}
  \bigl(a_{j,k,m}-q_{k,m}u_{j,k}\bigr)x_j,
&& k\in\mathcal{K},\,m\in\mathcal{M}_k,\\
& \sum_{j\in\mathcal{J}_z}x_j\le K_z,\qquad
  \sum_{j\in\mathcal{J}_z}s_jx_j\le P_z^{\mathrm{add}},\\
& \sum_{j\in\mathcal{J}_z(h)}x_j\le 1,
&& h\in\mathcal{H}_z,\\
& x_j\in\{0,1\},\quad e_c\ge 0,\quad d_{k,m}\ge 0.
\end{aligned}
\label{eq:mip_joint}
\end{equation}
Here $w_{k,m}=1/(|\mathcal{M}_k|\max(E_{z,k,m},1))$ gives every control table
equal weight and prevents large person tables from dominating; the
implementation divides $w$ by its maximum, a positive rescaling that does not
change the argmin. The released defaults are $\lambda_0=10^5$ and
$\lambda_1=10^{-2}$. With the default 256-person budget the whole marginal term
is below about $2.6\times10^4$, so covering one zero cell is worth more than
any marginal improvement the added set can buy. Coverage is still soft: an
unsourceable cell cannot make the program infeasible, and a cell can be left
uncovered when the budget or the candidate menu makes coverage impossible.
$K_z$ and $P_z^{\mathrm{add}}$ are small budgets (defaults: 64 households and
256 persons per zone), not the zone's population. The constraint over
$\mathcal{J}_z(h)\subseteq\mathcal{J}_z$ permits at most one transplant per
shell household.

The important difference from a per-cell objective is that
\eqref{eq:mip_joint} uses \emph{one} selection vector $x$ for \emph{all} control
tables. The deviations $d_{k,m}$ are therefore not independent cell repairs:
they measure how far the selected agents' aggregate profile is from the zone's
total marginal composition, table by table, and the solver must trade those
deviations against one another and against coverage. A set of records that fills
every zero cell but distorts the income or tenure shares is penalised; a set that
matches the shares but leaves a zero cell uncovered is penalised more heavily
through $\lambda_0$. This is a joint fit over the marginal tables, not a fit to
the full joint distribution: one selection vector must satisfy every table
simultaneously, while the joint distribution over attribute combinations is the
object of the zone fitter, not of the small pool-expansion program. The
coverage constraints are soft ($e_c$) precisely so that an unsourceable
cell---a structural zero, a padding code, or a cell with no donor anywhere---does
not make the program infeasible; it is reported instead.

The model is deliberately \emph{not} required to close the measured shortfall.
The deficit in a zone is often tens of thousands of persons, while the candidate
pool is a 5\% microsample and the added set is a few dozen households. Requiring
$\sum_j a_{j,k,m}x_j\ge E_{z,k,m}$ for every cell would either be infeasible or
force the pool to replicate the census many times over, which is the fitter's
job, not the pool's. What the expansion asks for is weaker and sufficient:
\emph{every zero cell for which a candidate exists should receive at least one
feasible candidate}, and the added candidates should jointly resemble the zone's
marginal totals. Because coverage is soft, a cell can still be left uncovered
when the budget or the candidate menu makes coverage impossible; such cells are
reported rather than silently dropped. The zone fit that follows decides which
candidates are actually used.

The module reports, per zone, the number of zero-capacity cells, how many are
covered, the selected household and person counts, the weighted joint-marginal
deviation, and the uncovered cells; the candidate-skip reasons are aggregated
over the run.

\emph{Development audit.} A ten-zone audit at the released defaults found 123
zero-capacity cells among 146 deficient cells (the other 23 were partial
shortfalls left to the zone fitter). The module generated donor--shell
candidates for 70 of the 73 sourceable zero cells and the joint selection
covered all 70, using 14 households (18 person rows) in total---one or two
households and at most three persons per zone, far inside the 64-household and
256-person budgets. All ten selection programs solved to optimality, and every
selected row passed the 11 shared rule checks. The 53 cells
left uncovered were 40 cells in four control codes absent from the census
extract, 10 \texttt{tend:6} padding cells, and three \texttt{tenlld:11} cases where no
donor--shell candidate survived the household and rule checks. The weighted
joint-marginal deviation of the added set was 0.239--1.622 per zone (median
0.406), and its unweighted mean absolute cell deviation was 0.126--0.280. These
are deviations of the added set, not of the fitted zone population; the audit
establishes that the mechanism can fill the sourceable zero cells with a small
number of feasible records, not that it closes the zone's full deficit.
Appendix~\ref{sec:pool_expansion_audit} gives the per-zone table and the
interpretation of the residual cells.
%% Audit source: outputs/full10_joint_report.json and
%% outputs/full10_joint_audit.log (10 zones, joint mode, released defaults;
%% 1,013.5 s).
